\chapter{第一讲补编}

重要的非线性波动方程, 如 Yang--Mills 方程、波映射与 Einstein 方程, 其基本对称性都与其几何特征紧密相联. 在研习这些方程时, 微分几何(特别是 Lorentz 几何)的知识, 与调和分析至少同等重要. 下面我们将考察 Minkowski 时空的共形结构与该结构同波动方程 (1) 之估计之间的关系. 特别地, 我们将用到第三讲中讨论的广义能量估计. 本文假定读者熟悉微分几何的基本概念, 如 Lie 微分、协变微分以及曲率张量.
\section*{4.1 Minkowski 时空的共形结构}

\setcounter{mtthm}{1}% 4.1 为节标题, 首个定理环境显示 4.2

\begin{definition}\label{def:4.2}
考虑一个时空 $(\mathcal{M}, g)$, 即一个 Lorentz 流形, 其中 $g$ 为号差 $(-1,1,\ldots,1)$ 的 Lorentz 度量.
一个微分同胚 $\Phi:\mathcal{U}\subset\mathcal{M}\to\mathcal{M}$ 称为\textbf{共形等距}, 如果对任意点 $p$ 均有 $\Phi_* g=\Lambda^2 g$, 即
\begin{equation}
(\Phi^* g)(X,Y)|_p=g(\Phi_* X,\Phi_* Y)|_{\Phi(p)}=\Lambda^2 g(X,Y)|_p
\end{equation}
且 $\Lambda\neq 0$. 若 $\Lambda=1$, 则称 $\Phi$ 为 $\mathcal{M}$ 的\textbf{等距}.\end{definition}

\begin{definition}\label{def:4.3}
若一个向量场 $K$ 生成一参数等距群(相应地, 共形等距群), 则称 $K$ 为\textbf{Killing 向量场}(相应地, \textbf{共形 Killing 向量场}).\end{definition}

设 $K$ 为这样的向量场, $\Phi_t$ 为相应的单参数群. 由于 $(\Phi_t)_*$ 均为共形等距, 可知 $\mathcal{L}_K g$ 必与度量 $g$ 成比例; 若 $K$ 为 Killing 向量场, 则 $\mathcal{L}_K g=0$.
注意, C(共形结构)关于共形等距具有共形不变性这一重要性质.
\begin{definition}\label{def:4.4}
对任意向量场 $X$, 记 $(X)_\pi$ 为 $X$ 的\textbf{形变张量}, 由下式定义:
\begin{equation}
{}^{(X)}\pi_{\alpha\beta}=(\mathcal{L}_X g)_{\alpha\beta}=D_\alpha X_\beta+D_\beta X_\alpha.\end{equation}
\end{definition}

该张量 $(X)\pi$ 以一种精确的意义, 衡量了由 $X$ 生成的微分同胚偏离等距或共形等距的程度.
\begin{proposition}\label{prop:4.5}
向量场 $X$ 为 Killing 向量场, 当且仅当 $(X)\pi=0$; 为共形 Killing 向量场, 当且仅当 $(X)\pi$ 与 $g$ 成比例.\end{proposition}

为证明"当"的部分, 可取局部坐标 $x^0,x^1,\ldots,x^n$ 使得 $X=\frac{\partial}{\partial x^0}$. 由此立即可见, 相对这些坐标, 度量 $g$ 与 $x^0$ 无关.
\begin{proposition}\label{prop:4.6}
任意时空 $\mathcal{M}$ 上, 线性无关的 Killing 向量场个数至多为 $\frac12(n+1)(n+2)$.\end{proposition}

拥有最多 Killing 与共形 Killing 向量场的时空正是 Minkowski 时空. 下面我们复习其相应的等距与共形等距.
设 $x^\mu$ 为正定向的惯性坐标系, 我们有:
\begin{enumerate}
\item 平移: 对任意给定向量 $a=(a^0,a^1,\ldots,a^n)\in\mathbf{M}$,
\begin{equation}
x^\mu\to x^\mu+a^\mu.\end{equation}
\item Lorentz 旋转: 对任意 $\Lambda=\Lambda^\rho_\sigma\in\mathbf{O}(1,n)$,
\begin{equation}
x^\mu\to \Lambda^\mu_\nu x^\nu.\end{equation}
\item 伸缩: 对任意实数 $\lambda\neq 0$,
\begin{equation}
x^\mu\to \lambda x^\mu.\end{equation}
\item 反演: 考虑变换 $x^\mu\to I(x^\mu)$, 其中
\begin{equation}
I(x^\mu)=\frac{x^\mu}{(x,x)},\end{equation}
对一切满足 $(x,x)\neq 0$ 的点 $x\in\mathbf{M}$ 定义.\end{enumerate}

前两类变换是 $\mathbf{M}$ 的等距, 由它们生成的群称为 Poincar\'e 群; 后两类变换是共形等距, 由上述全部变换生成的群称为共形群. 下面列出生成上述变换的 Killing 与共形 Killing 向量场.
\begin{enumerate}[label=\roman*.]
\item $x^\mu$ 方向的平移生成元, $\mu=0,1,\ldots,n$:
\begin{equation}
T_\mu=\frac{\partial}{\partial x^\mu}.\end{equation}
\item $(\mu,\nu)$ 平面中 Lorentz 旋转的生成元:
\begin{equation}
L_{\mu\nu}=x_\mu\partial_\nu-x_\nu\partial_\mu.\end{equation}
\item 伸缩变换的生成元:
\begin{equation}
S=x^\mu\partial_\mu.\end{equation}
\item 反演平移的生成元:
\begin{equation}
K_\mu=2x_\mu\Bigl(x^\rho\frac{\partial}{\partial x^\rho}-(x^\rho x_\rho)\frac{\partial}{\partial x^\mu}\Bigr).\end{equation}
\end{enumerate}

\begin{theorem}\label{thm:4.7}
$\mathbf{M}^{n+1}$ 中任意共形 Killing 向量场($n>1$)都是向量场 $T,L,S,K$ 以实常系数为系数的线性组合. 若 $n=1$, 则所有共形向量场由下式给出:
\begin{equation}
f(t+x)(\partial_t+\partial_x)+g(t-x)(\partial_t-\partial_x),\end{equation}
其中 $f,g$ 为任意函数.\end{theorem}

\setcounter{mtthm}{8}% 4.8 为节标题, 下一定理环境显示 4.9

\section*{4.8 广义能量估计}

考虑平直 Minkowski 空间 $\mathbb{R}^{n+1}$ 中的波动方程 $\Box\phi=0$ 及其相伴的能量---动量张量
\begin{equation}
Q_{\alpha\beta}=Q_{\alpha\beta}[\phi]=\partial_\alpha\phi\partial_\beta\phi-\frac12 m_{\alpha\beta}\bigl(m^{\mu\nu}\partial_\mu\phi\partial_\nu\phi\bigr).\end{equation}

\begin{proposition}\label{prop:4.9}
能量---动量张量是对称且无散度的, 即
\begin{equation}
\partial^\beta Q_{\alpha\beta}=0.\end{equation}
此外, 对任意指向未来的类时向量场 $X,Y$, 均有 $Q(X,Y)\geq 0$.\end{proposition}

\begin{proof}
前两个结论显然. 为验证第三个, 在任一点取一对零向量 $L,\underline{L}$ 张成包含 $X,Y$ 的平面, 于是 $X,Y$ 都是 $L,\underline{L}$ 带正系数的线性组合. 由如下简单而重要的计算可知其正性:
\begin{gather}
Q(L,L)=L(\phi)^2,\\
Q(L,\underline{L})=|\bar\nabla\phi|^2,\\
Q(\underline{L},\underline{L})=L(\phi)^2,\end{gather}
其中 $|\bar\nabla\phi|=\sum_A |e_A(\phi)|^2$, $(e_A)_{A=1,\ldots,n-1}$ 为张成 $L,\underline{L}$ 正交补的标准正交标架.\end{proof}

设 $X$ 为任意向量场, 考虑与之相伴的时空动量向量场 $P_\alpha=Q_{\alpha\beta}X^\beta$. 利用 $Q$ 的对称性与无散度性质, 连同形变张量 $\pi=(X)\pi$ 的定义, 推得
\begin{equation}
\partial_\alpha P^\alpha=\frac12 Q^{\alpha\beta}\pi_{\alpha\beta}.\end{equation}

因此, 若 $X$ 为 Killing 向量场, 则 $\partial_\alpha P^\alpha=0$, 积分后即得熟知的守恒律. 特别地, 若我们在时空板 $[t_0,t]\times\mathbb{R}^n$ 上积分, 便得到
\begin{equation}
\int_{\Sigma_t}Q(X,\partial_t)\,dx=\int_{\Sigma_{t_0}}Q(X,\partial_t)\,dx,\end{equation}
其中 $\partial_t$ 在此扮演类空超曲面 $\Sigma_t$ 的未来指向单位法向. 取 $X$ 亦为 Killing 向量场 $\partial_t$, 即得熟知的能量守恒律
\begin{equation}
\int_{\Sigma_t}Q(\partial_t,\partial_t)\,dx=\int_{\Sigma_{t_0}}Q(\partial_t,\partial_t)\,dx,\end{equation}
其中 $Q(\partial_t,\partial_t)=\frac12\bigl(|\partial_t\phi|^2+\sum_{i=1}^n|\partial_i\phi|^2\bigr)$ 为通常的能量密度. 线性动量与角动量的守恒律也包含在上式中, 只需取 $X=\partial_i$ 或 $X=L_{ij}$.
若我们不在板 $[0,T]\times\mathbb{R}^n$ 上积分恒等式 $\partial^\alpha P_\alpha$, 而是在未来区域
\begin{equation}
\mathcal{D}=\{(t,x)\mid |x-y|\leq T-t,\; t_0\leq t\leq t_1\leq T\}
\end{equation}
上积分, 则得到守恒律
\begin{equation}
\int_{\Sigma_{t_1}}Q(X,\partial_t)\,dx+\int_{\underline{C}(t_0,t_1)}Q(X,\underline{L})=\int_{\Sigma_{t_0}}Q(X,\partial_t)\,dx,\end{equation}
其中 $\underline{C}(t_0,t_1)$ 为 $\mathcal{D}$ 的外边界, $\underline{L}=\partial_t-\sum_i\frac{(x^i-x_0^i)}{|x-x_0|}\partial_i$ 为指向入射零超曲面 $\underline{C}$ 的法向零向量场, 而通量 $\int_{\underline{C}(t_0,t_1)}Q(X,\underline{L})=\int_{t_0}^{t_1}dt\int_{S_{\underline{u}_t}}Q(X,\underline{L})$, 其中 $S_{\underline{u}_t}$ 为 $\underline{C}$ 与 $\Sigma_t$ 的交球面. 在 $X=\partial_t$ 的特殊情形, $ \Sigma_t $ 上的能量积分与通量均为正, 从而给出波动方程依赖域性质的简单证明.
上面导出 (97) 的程序可推广到共形 Killing 向量场.
\begin{proposition}\label{prop:4.10}
设 $Q_{\alpha\beta}=Q_{\alpha\beta}[\phi]$ 为 $\Box\phi=0$ 之解相伴的能量---动量张量. 令 $X$ 为共形 Killing 向量场, 即 $\pi=(X)\pi=\mathcal{L}_X m=\Omega m$, 且 $\mathrm{tr}\,\pi=m^{\alpha\beta}\pi_{\alpha\beta}$. 易验证 $\Box\Omega=0$; 特别地, 当 $X=K_0$ 时 $\Omega=4(n+1)t$. 记
\begin{equation}
\bar P_\alpha=Q_{\alpha\beta}X^\beta+\frac{n-1}{4(n+1)}\mathrm{tr}\,\pi\,\phi\partial_\alpha\phi-\frac{n-1}{8(n+1)}\partial_\alpha(\mathrm{tr}\,\pi)\phi^2.\end{equation}
则有
\begin{equation}
\partial^\alpha\bar P_\alpha=0.\end{equation}
\end{proposition}

将命题应用于共形 Killing 向量场 $X=K:=K_0$, 并在时间板 $[t_0,t]\times\mathbb{R}^n$ 上积分相应的无散度方程, 推得:
\begin{proposition}\label{prop:4.11}
设 $\overline{Q}(K_0,T_0)=Q(K_0,T_0)+(n-1)t\phi\partial_t\phi-\frac{n-1}{2}\phi^2$, 其中 $T_0=\partial_t$ 为 $\Sigma_t$ 的单位法向, $\phi$ 为 $\Box\phi=0$ 的解.
\begin{enumerate}[label=\roman*.]
\item 如下共形守恒律成立:
\begin{equation}
\int_{\Sigma_t}\overline{Q}(K_0,T_0)=\int_{\Sigma_{t_0}}\overline{Q}(K_0,T_0).\end{equation}
\item 此外,
\begin{equation}
\int_{\Sigma_t}\overline{Q}(K_0,T_0)=\frac14\Bigl(\int_{\Sigma_t}\underline{u}^2(L'\phi)^2+\int_{\Sigma_t}2(t^2+r^2)|\nabla\phi|^2+\int_{\Sigma_t}u^2(\underline{L}'\phi)^2\Bigr),\end{equation}
其中 $L=\partial_t+\partial_r$, $\underline{L}=\partial_t-\partial_r$, $u=t-r$, $\underline{u}=t+r$, 且 $\underline{u}L'(\phi)=\underline{u}L(\phi)+(n-1)\phi$, $u\underline{L}'(\phi)=u\underline{L}(\phi)+(n-1)\phi$.
\item 又若 $n\geq 3$, 则存在常数 $c>0$ 使得
\begin{equation}
\int_{\Sigma_t}\overline{Q}(K_0,T_0)\geq c\Bigl(\int_{\Sigma_t}\underline{u}^2(L\phi)^2+\int_{\Sigma_t}2(t^2+r^2)|\nabla\phi|^2+\int_{\Sigma_t}u^2(\underline{L}\phi)^2+\phi^2\Bigr).\end{equation}
\end{enumerate}
\end{proposition}

\begin{definition}\label{def:4.12}
我们引入如下范数:
\begin{gather}
E[\phi](t)=\int_{\Sigma_t}|\phi(t)|^2,\\
\mathcal{Q}[\phi](t)=\int_{\Sigma_t}\overline{Q}[\phi](K_0,T_0),\\
\mathcal{E}[\phi](t)=\int_{\Sigma_t}\Bigl(\underline{u}^2|L\phi|^2+2(t^2+r^2)|\nabla\phi|^2+u^2|\underline{L}\phi|^2+|\phi|^2\Bigr).\end{gather}
我们还引入高阶范数:
\begin{gather}
E_s[\phi](t)=\sum_{j=0}^s\;\sum_{\Gamma_{i_1},\ldots,\Gamma_{i_j}}E[\Gamma_{i_1},\ldots,\Gamma_{i_j}\phi](t),\\
\mathcal{Q}_s[\phi](t)=\sum_{j=0}^s\;\sum_{\Gamma_{i_1},\ldots,\Gamma_{i_j}}\mathcal{Q}[\Gamma_{i_1},\ldots,\Gamma_{i_j}\phi](t),\end{gather}
其中 $\Gamma$ 为任一向量场
\begin{equation}
S=t\partial_t+\sum_i x_i\partial_i,\qquad L_i=x_i\partial_t+t\partial_i,\qquad \Omega_{ij}=x_i\partial_j-x_j\partial_i.\end{equation}
\end{definition}

\setcounter{mtthm}{12}% 4.13 与 4.14 之间无跳过, 但 4.13 为注, 需钉号

\begin{remark}\label{rem:4.13}
注意
\begin{equation}
E_{s+1}(t)\lesssim\mathcal{Q}_s(t)=\mathcal{Q}_s(0).\end{equation}
\end{remark}

\setcounter{mtthm}{14}% 4.14 为节标题, 下一定理环境显示 4.15

\section*{4.14 衰减估计}

众所周知, 自由波动方程 $\Box\phi=0$ 在 $\mathbb{R}^{n+1}$ 中的解, 在一致范数下随 $t\to\infty$ 按 $t^{-\frac{n-1}{2}}$ 衰减. 此外, 远离出射特征方向时衰减更快. 这类信息可由物理空间或傅里叶空间中解的显式表示导出. 本节的目标是说明: 无需使用显式表示, 仅借助前一节讨论的广义能量估计, 便能获得此类信息.
考虑标准的零向量对 $E_\pm=\partial_t\pm\partial_r$, 以及角向量场 $A_i=\partial_i-\frac{x_i}{r}\partial_r$. 显然
\begin{equation}
\sum_i|A_i\phi|\lesssim|\bar\nabla\phi|\lesssim\sum_i|A_i\phi|.\end{equation}
又有
\begin{equation}
|\partial_r\phi|+\sum_i|A_i\phi|\lesssim|D\phi|\lesssim|\partial_r\phi|+\sum_i|A_i\phi|.\end{equation}
我们还可轻易验证如下简单的代数恒等式:
\begin{gather}
\frac12\bigl((t+r)E_+ +(t-r)E_-\bigr)=S,\\
\frac12\bigl((t+r)E_+ -(t-r)E_-\bigr)=\sum_i\frac{x_i}{|x|}L_i,\\
tA_i=L_i-\frac{x_i}{|x|}\sum_j\frac{x_j}{|x|}L_j,\\
t\Omega_{ij}=x_iL_j-x_jL_i.\end{gather}
由前两个恒等式易得
\begin{gather}
|E_+\phi(t,x)|\lesssim\frac1t|\Gamma\phi(t,x),\\
|E_-\phi(t,x)|\lesssim\frac1{|t-|x||}|\Gamma\phi(t,x)|,\end{gather}
其中 $|\Gamma\phi|=|S\phi|+|L\phi|$. 以及
\begin{equation}
|\bar\nabla\phi(t,x)|\lesssim\frac1t|\Gamma\phi(t,x)|.\end{equation}
显然还有
\begin{equation}
|\partial\phi(t,x)|\lesssim\frac1{|t-|x||}|\Gamma\phi(t,x)|,\end{equation}
或更一般地
\begin{equation}
|\partial^N\phi(t,x)|\lesssim\frac1{|t-|x||^N}|\Gamma^N\phi(t,x)|,\end{equation}
其中 $|\Gamma^N\phi|=\sum|\Gamma_1\ldots\Gamma_N\phi|$, $\Gamma_1,\ldots,\Gamma_N$ 为任意向量场 $S,L_1,\ldots,L_n$.
将上述不等式与我们所定义的范数结合, 推得
\begin{gather}
t\|E_+\phi(t)\|_{L^2}\lesssim E_1(\phi)^{\frac12},\\
t\|\nabla\phi(t)\|_{L^2}\lesssim E_1(\phi)^{\frac12},\\
\|uE_-\phi(t)\|_{L^2}\lesssim E_1(\phi)^{\frac12}.\end{gather}

现在我们准备证明如下结论:
\begin{proposition}\label{prop:4.15}
设 $\Box\phi=0$, 初数据满足上述假设. 则对所有 $t\geq\Lambda_0$, $s>\frac n2$,
\begin{gather}
\|\phi(t)\|_{L^\infty}\lesssim\Bigl(\frac1t\Bigr)^{\frac{n-1}{2}}E_s[\phi](t)^{\frac12},\\
\|(1+|u|)^k\partial^k\phi(t)\|_{L^\infty}\lesssim\Bigl(\frac1t\Bigr)^{\frac{n-1}{2}}E_{s+k}[\phi](t)^{\frac12}.\end{gather}
此外
\begin{gather}
\|E_+\phi(t)\|_{L^\infty}\lesssim\Bigl(\frac1t\Bigr)^{\frac{n+1}{2}}E_{s+1}[\phi](t)^{\frac12},\\
\|\bar\nabla\phi(t)\|_{L^\infty}\lesssim\Bigl(\frac1t\Bigr)^{\frac{n+1}{2}}E_{s+1}[\phi](t)^{\frac12},\\
\|(1+|u|)E_-\phi(t)\|_{L^\infty}\lesssim\Bigl(\frac1t\Bigr)^{\frac{n-1}{2}}E_{s+1}[\phi](t)^{\frac12}.\end{gather}
\end{proposition}

证明基于如下引理:
\begin{lemma}\label{lem:4.16}
设 $u(x)$ 为 $\mathbb{R}^n$($n\geq 2$)上光滑、具紧支集的函数. 则
\begin{equation}
|u(x)|\leq C\frac1{|x|^{n-1}}\Bigl(\|\partial_r u\|_{L^1}+\|(r\bar\nabla)^{n-1}\partial_r u\|_{L^1}\Bigr).\end{equation}
\end{lemma}

作为该引理的直接推论, 我们得到
\begin{equation}
|u(x)|\leq C\frac1{|x|^{n-1}}\sum_{k=0}^{n-1}\|\partial_r\Omega^k u\|_{L^1},\end{equation}
其中 $\Omega^k u$ 表示 $\Omega_1\ldots\Omega_k u$ 形式的所有导数, $\Omega_i=\epsilon_{ijk}x_j\partial_k$ 为角动量向量场. 现对固定某 $t$ 的函数 $u(x)=\phi^2(t,x)$ 应用不等式 (108):
\begin{equation}
|\phi(t,x)|^2\lesssim\frac1{|x|^{n-1}}\sum_{k=0}^{n-1}\sum_{i+j=k}\|\Omega^i\phi(t)\|_{L^2}\|\partial\Omega^j\phi(t)\|_{L^2}\lesssim\frac1{|x|^{n-1}}E_n[\phi].\end{equation}
由此推知, 当 $|x|\geq\frac t2$ 时,
\begin{equation}
|\phi(t,x)|\lesssim\Bigl(\frac1t\Bigr)^{\frac{n-1}{2}}E_n^{\frac12}[\phi].\end{equation}

现考虑区域 $|x|\leq\frac t2$. 取支在 $|x|\leq 1$ 内的检验函数 $\chi$, 且当 $|x|\leq\frac12$ 时 $\chi(x)=1$. 对函数 $u(x)=\chi(x/t)\phi^2(t,x)$ 应用 Sobolev 不等式
\begin{equation}
|u(x)|\lesssim\|D^n u\|_{L^1}.\end{equation}
于是当 $|x|\leq\frac t2$ 时, 顾及 (104) 再顾及 (31),
\begin{align*}
|\phi(t,x)|^2&\lesssim\sum_{i+j+k=n}\|D^i\chi(x/t)\|_{L^\infty}\cdot\|D^j\phi(t)\|_{L^2(|x|\leq\frac t2)}\cdot\|D^k\phi(t)\|_{L^2(|x|\leq\frac t2)}\\
&\lesssim\sum_{i+j+k=n}t^{-i}\cdot t^{-j}\|\Gamma^j\phi(t)\|_{L^2}\cdot t^{-k}\|\Gamma^k\phi(t)\|_{L^2}\\
&\lesssim t^{-n}E_n[\phi].\end{align*}
因此
\begin{equation}
\sup_{|x|\leq\frac t2}|\phi(t,x)|\lesssim\Bigl(\frac1t\Bigr)^{\frac n2}E_n^{\frac12}[\phi],\end{equation}
这与 (109) 一起证明了
\begin{equation}
\|\phi(t)\|_{L^\infty}\lesssim\Bigl(\frac1t\Bigr)^{\frac{n-1}{2}}E^{\frac12}[\phi],\end{equation}
正如所愿.
为证明\lemref{lem:4.16}, 在极坐标下写 $x=r\xi$, $\xi\in\mathbb{R}^{n-1}$:
\begin{equation}
u(r\xi)=-\int_r^\infty\partial_r u(\lambda\xi)\,d\lambda.\end{equation}
于是
\begin{equation}
\int_{|\xi|=1}|u(r\xi)|\,d\sigma(\xi)\leq\frac1{r^{n-1}}\int_{\mathbb{R}^n}|\partial_r u(y)|\,dy.\end{equation}
另一方面, 在单位球面 $S^{n-1}$ 上使用 Sobolev 不等式,
\begin{equation}
|u(x)|\leq c_n\Bigl(\|u(r\cdot)\|_{L^1(S^{n-1})}+\|(r\bar\nabla)^{n-1}u(r\cdot)\|_{L^1(S^{n-1})}\Bigr),\end{equation}
结合上面的不等式即证得所需结果.
\begin{definition}\label{def:4.17}
设 $\phi,\psi$ 为 $\mathbb{R}^{n+1}$ 上两个光滑函数. 我们定义
\begin{gather}
Q_0(\phi,\psi)=m^{\alpha\beta}\partial_\alpha\phi\partial_\beta\psi,\\
Q_{\alpha\beta}(\phi,\psi)=\partial_\alpha\phi\partial_\beta\psi-\partial_\alpha\psi\partial_\beta\phi.\end{gather}
\end{definition}

下面我们将利用前面建立的衰减估计, 并仿照 \cite{klai28} 的精神对零形式作适当分解, 来估计零二次型 $Q_0(\phi,\psi)$ 与 $Q_{ij}(\phi,\psi)$. 设 $\Box\phi=\Box\psi=0$. 考虑零形式 $Q_0(\phi,\psi)$ 并将其写成
\begin{equation}
Q_0(\phi,\psi)=-\frac12\bigl(E_+\phi\cdot E_-\psi+E_+\phi\cdot E_-\psi\bigr)+\bar\nabla\phi\cdot\bar\nabla\psi.\end{equation}
于是由\propref{prop:4.15} 的结果,
\begin{align*}
\|Q_0(\phi,\psi)(t)\|_{L^2}&\leq\frac12\Bigl(\|E_+(\phi)(t)\|_{L^\infty}\|E_-(\psi)(t)\|_{L^2}+\|E_-(\phi)(t)\|_{L^\infty}\|E_+(\psi)(t)\|_{L^2}\Bigr)\\
&\quad+\|\nabla\phi(t)\|_{L^\infty}\|\nabla\psi(t)\|_{L^2}\\
&\lesssim\Bigl(\frac1t\Bigr)^{\frac{n+1}{2}}E_{n+1}[\phi]^{\frac12}E_1[\psi]^{\frac12}.\end{align*}
类似地, 为估计 $Q_{ij}(\phi,\psi)$, 写 $\partial_i=A_i+\frac{x_i}{|x|}\partial_r$. 于是
\begin{equation}
Q_{ij}(\phi,\psi)=A_i\phi A_j\psi-A_j\phi A_i\psi+\partial_r\phi\Bigl(\frac{x_i}{r}A_j\phi-\frac{x_j}{r}A_i\Bigr)\phi.\end{equation}
顾及 (32) 中第三个公式以及 (31), 我们对 $Q_{ij}$ 得到同类型的估计.
